\documentclass[11pt]{amsart}

\usepackage[T1]{fontenc}
\usepackage{lmodern}
\usepackage{microtype}
\usepackage{mathtools,amssymb,amsthm}
\usepackage{array}
\usepackage[hidelinks]{hyperref}

\hypersetup{
  pdftitle={An Infinite Family of Polyominoes Whose C4-Face-Magic Spectrum Omits Its Centre}
}

\emergencystretch=2em

\newtheorem{theorem}{Theorem}
\newtheorem{lemma}[theorem]{Lemma}

\DeclareMathOperator{\Spec}{Spec}
\newcommand{\cells}{\mathcal C}

\title[A non-interval $C_4$-face-magic spectrum]
{An Infinite Family of Polyominoes Whose\\ $C_4$-Face-Magic Spectrum Omits Its Centre}

\date{}

\subjclass[2020]{05C78, 05B50, 05C10}
\keywords{$C_4$-face-magic labeling, polyomino, magic spectrum, face-magic
graph, counterexample}

\begin{document}

\begin{abstract}
For a polyomino $P$ with $n=|V(P)|$, Chalise, Low and
Eisenkolb-Vaithyanathan define the $C_4$-face-magic spectrum
$\Spec_{C_4}(P)$ to be the set of magic constants of $C_4$-face-magic
labelings of $P$, and ask (Question~2 of their Section~5) whether
$\Spec_{C_4}(P)$ is always an interval of integers. It is not. For the
L-tromino $L_2$ we determine the spectrum exactly,
\[
  \Spec_{C_4}(L_2)=\{15,16,17,19,20,21\},
\]
so the centre $2(n+1)=18$ is missing. The mechanism is a double count: if
some multiset of cells covers $V(P)$ with multiplicity one except for a
single uncovered vertex and a single doubled vertex, then $4\mid n$, the
value $2(n+1)$ is excluded, and $\Spec_{C_4}(P)$ is confined to the six
values $2(n+1)+j$ with $1\le|j|\le3$. An explicit L-shaped family $L_m$,
$m$ even, satisfies the hypothesis for every $m$, so the centre is omitted
at every order $n=2m+4$ with $m$ even. Both the exclusion and the six-value
window are proved without any search.
\end{abstract}

\maketitle

\section{Introduction}

A \emph{polyomino} $P$ is a finite edge-connected union of unit squares in the
square lattice $\mathbb Z^2$, identified with its associated plane graph, whose
vertices are the lattice vertices of the cells and whose edges are the lattice
edges \cite[Section~1]{CLE}. We write $\cells(P)$ for the set of cells,
$V(P)$ for the vertex set, $n=|V(P)|$ and $[n]=\{1,\ldots,n\}$. Throughout,
the cell $[x,y]$ is the closed unit square with lower-left lattice point
$(x,y)$, so
\[
  V\bigl([x,y]\bigr)=\bigl\{(x,y),\,(x+1,y),\,(x,y+1),\,(x+1,y+1)\bigr\}.
\]

Chalise, Low and Eisenkolb-Vaithyanathan \cite{CLE} prove the Shiu--Low--Liu
conjecture \cite{SLL} that every polyomino is $C_4$-face-magic, and in their
Section~5 they introduce the \emph{$C_4$-face-magic spectrum}
\begin{equation}\label{eq:spec}
\begin{gathered}
  \Spec_{C_4}(P)
  =\Bigl\{
     s\in\mathbb N:\ \exists\ \text{a bijection }\lambda:V(P)\to[n]\\
     \text{such that }\textstyle\sum_{v\in V(C)}\lambda(v)=s
     \text{ for every }C\in\cells(P)
   \Bigr\}.
\end{gathered}
\end{equation}
They observe that $\Spec_{C_4}(P)$ is nonempty for every polyomino, that the
complementary labeling $\overline\lambda=n+1-\lambda$ gives the symmetry
\begin{equation}\label{eq:sym}
  s\in\Spec_{C_4}(P)
  \iff
  4(n+1)-s\in\Spec_{C_4}(P),
\end{equation}
and that $\Spec_{C_4}(P)\subseteq\{10,11,\ldots,4n-6\}$. Because of
\eqref{eq:sym} the point $2(n+1)$ is the \emph{centre} of the spectrum. They
then pose five questions, the second of which reads, verbatim,
\begin{quote}
  ``Is $\Spec_{C_4}(P)$ always an interval of integers?''
\end{quote}
\cite[Section~5, Question~2]{CLE}. We answer it negatively; a witness has
three cells, and the obstruction is a counting identity that needs no
computer.

\begin{theorem}\label{thm:tromino}
Let $L_2$ be the L-tromino, the polyomino with cells $[0,0]$, $[0,1]$,
$[1,0]$. Then $n=|V(L_2)|=8$, the centre is $2(n+1)=18$, and
\[
  \Spec_{C_4}(L_2)=\{15,16,17,19,20,21\}.
\]
In particular $18\notin\Spec_{C_4}(L_2)$ while $15$ and $21$ do lie in it, so
$\Spec_{C_4}(L_2)$ is not an interval of integers and Question~2 of
\cite{CLE} is false.
\end{theorem}

\begin{theorem}\label{thm:family}
For even $m\ge2$ let $L_m$ be the polyomino with cells
$[0,0],[0,1],\ldots,[0,m-1]$ together with $[1,0]$, so that $L_2$ is the
L-tromino and $L_4$ the L-pentomino. Then $n=|V(L_m)|=2m+4$,
\[
\begin{gathered}
  2(n+1)=4m+10\notin\Spec_{C_4}(L_m),\\
  \Spec_{C_4}(L_m)\subseteq\bigl\{\,2(n+1)+j:\ 1\le|j|\le3\,\bigr\}.
\end{gathered}
\]
Since $\Spec_{C_4}(L_m)$ is nonempty and symmetric about $2(n+1)$, it
contains values on both sides of the missing centre, so it is not an interval.
Question~2 of \cite{CLE} therefore fails for infinitely many polyominoes, at
infinitely many orders.
\end{theorem}

Both mechanisms used below are in print and are references of \cite{CLE}:
Curran \cite[p.~362]{Curran} confines the $C_4$-face-magic value of an
odd-order projective grid graph $P_{m,n}$ to $\{2mn+1,2mn+2,2mn+3\}$ by the
same kind of double count, and the label-collision device is Lemma~2.1 of
Shiu--Low--Liu \cite{SLL}.

\section{Two lemmas}\label{sec:lemmas}

\begin{lemma}\label{lem:cover}
Let $P$ be a polyomino with $n=|V(P)|$, and suppose there is a finite multiset
$S$ of cells of $P$ and two distinct vertices $p,v\in V(P)$ such that, writing
$w(u)$ for the number of members of $S$ whose vertex set contains $u$,
\[
  w(p)=0,\qquad w(v)=2,\qquad w(u)=1\ \text{ for all }u\in V(P)\setminus\{p,v\}.
\]
Call such an $S$ a \emph{single-defect cover} of $P$. Then
\begin{enumerate}
\item[(i)] $4\mid n$ and $|S|=n/4$;
\item[(ii)] every $s\in\Spec_{C_4}(P)$ satisfies
  $\;s=2(n+1)+\dfrac{4\delta}{n}\;$ for some integer $\delta$ with
  $1\le|\delta|\le n-1$;
\item[(iii)] $2(n+1)\notin\Spec_{C_4}(P)$.
\end{enumerate}
\end{lemma}

\begin{proof}
Each cell has exactly four vertices, so
$\sum_{u\in V(P)}w(u)=4|S|$. On the other hand the hypothesis gives
$\sum_{u}w(u)=0+2+(n-2)=n$. Hence $4|S|=n$, which is (i).

Let $\lambda$ be a $C_4$-face-magic labeling of $P$ with magic constant $s$,
and put $N=\sum_{i=1}^{n}i=n(n+1)/2$. Summing the magic condition over the
members of $S$ and exchanging the order of summation,
\[
  |S|\,s
  =\sum_{C\in S}\ \sum_{u\in V(C)}\lambda(u)
  =\sum_{u\in V(P)}w(u)\,\lambda(u)
  =N-\lambda(p)+\lambda(v).
\]
Writing $\delta=\lambda(v)-\lambda(p)$ and $|S|=n/4$ this is
$(n/4)s=N+\delta$, i.e.
\begin{equation}\label{eq:key}
  s=\frac{4(N+\delta)}{n}=2(n+1)+\frac{4\delta}{n}.
\end{equation}
Since $\lambda$ is a bijection onto $[n]$ and $p\ne v$, we have
$\lambda(p)\ne\lambda(v)$, so $\delta\ne0$, and both labels lie in $[n]$, so
$|\delta|\le n-1$. That is (ii).

For (iii), take $s=2(n+1)$ in \eqref{eq:key}: it forces $\delta=0$, i.e.
$\lambda(p)=\lambda(v)$ with $p\ne v$, contradicting injectivity. Equivalently
and with no algebra: $(n/4)\cdot2(n+1)=n(n+1)/2=N$, so the identity reads
$N=N-\lambda(p)+\lambda(v)$.
\end{proof}

\begin{lemma}\label{lem:window}
Under the hypothesis of Lemma~\ref{lem:cover}, write $n=4t$. Then
\[
  \Spec_{C_4}(P)\subseteq\bigl\{\,2(n+1)+j:\ j\in\{-3,-2,-1,1,2,3\}\,\bigr\},
\]
so $|\Spec_{C_4}(P)|\le6$, however large $P$ is.
\end{lemma}

\begin{proof}
By \eqref{eq:key}, $s-2(n+1)=4\delta/n=\delta/t$. As $s$ is an integer,
$t\mid\delta$, say $\delta=jt$ with $j\in\mathbb Z$. By
Lemma~\ref{lem:cover}(ii), $\delta\ne0$ and $|\delta|\le n-1=4t-1$, so
$0<|j|\le4-1/t<4$ and hence $j\in\{\pm1,\pm2,\pm3\}$. Finally
$s=2(n+1)+j$.
\end{proof}

So a polyomino with a single-defect cover has a spectrum contained in the six
integers within distance $3$ of its centre, and missing the centre itself.
Nothing but a covering count has been used; in particular no labeling has been
searched for.

\section{The L-tromino}\label{sec:tromino}

\begin{proof}[Proof of Theorem~\ref{thm:tromino}]
The cells of $L_2$ are $[0,0],[0,1],[1,0]$, so
\[
  V(L_2)=\{(0,0),(0,1),(0,2),(1,0),(1,1),(1,2),(2,0),(2,1)\},
\]
$n=8$, $t=n/4=2$ and the centre is $2(n+1)=18$.

\emph{Upper bound.} Take $S=\{[0,1],\,[1,0]\}$. Its two cells have vertex sets
$\{(0,1),(1,1),(0,2),(1,2)\}$ and $\{(1,0),(2,0),(1,1),(2,1)\}$, so $(1,1)$ is
covered twice, $(0,0)$ is covered not at all, and each of the remaining five
vertices exactly once. Thus $S$ is a single-defect cover with $p=(0,0)$,
$v=(1,1)$ and $|S|=2=n/4$. By Lemmas~\ref{lem:cover} and \ref{lem:window},
\[
  \Spec_{C_4}(L_2)\subseteq\{15,16,17,19,20,21\}.
\]

\emph{Lower bound.} Each of the six columns of Table~\ref{tab:L2} is a
bijection $V(L_2)\to[8]$ under which all three cells of $L_2$ sum to the
value heading the column. Hence each of the six values lies in
$\Spec_{C_4}(L_2)$, and the two bounds meet.

Since $15,21\in\Spec_{C_4}(L_2)$ and $18\notin\Spec_{C_4}(L_2)$, the set is
not an interval.
\end{proof}

\begin{table}[ht]
\[
\begin{array}{c|cccccc}
 u & \lambda_{15}(u) & \lambda_{16}(u) & \lambda_{17}(u)
   & \lambda_{19}(u) & \lambda_{20}(u) & \lambda_{21}(u)\\
\hline
(0,0) & 7 & 5 & 3 & 1 & 1 & 1\\
(0,1) & 5 & 8 & 8 & 8 & 8 & 8\\
(0,2) & 3 & 3 & 2 & 2 & 3 & 2\\
(1,0) & 2 & 2 & 5 & 7 & 6 & 5\\
(1,1) & 1 & 1 & 1 & 3 & 5 & 7\\
(1,2) & 6 & 4 & 6 & 6 & 4 & 4\\
(2,0) & 4 & 6 & 4 & 4 & 2 & 3\\
(2,1) & 8 & 7 & 7 & 5 & 7 & 6
\end{array}
\]
\caption{Six $C_4$-face-magic labelings of the L-tromino $L_2$, one for each
value of $\Spec_{C_4}(L_2)=\{15,16,17,19,20,21\}$; $\lambda_s$ has magic
constant $s$. For example $\lambda_{17}$ gives cell $[0,0]$ the sum
$3+5+8+1=17$, cell $[0,1]$ the sum $8+1+2+6=17$ and cell $[1,0]$ the sum
$5+4+1+7=17$.}
\label{tab:L2}
\end{table}

\section{The infinite family}\label{sec:family}

\begin{proof}[Proof of Theorem~\ref{thm:family}]
Fix an even $m\ge2$. The $m$ column cells $[0,j]$, $0\le j\le m-1$, together
contribute the $2(m+1)$ vertices $\{0,1\}\times\{0,\ldots,m\}$, and the cell
$[1,0]$ contributes in addition the two vertices $(2,0)$ and $(2,1)$. Hence
\[
\begin{gathered}
  V(L_m)=\bigl(\{0,1\}\times\{0,\ldots,m\}\bigr)\cup\bigl\{(2,0),(2,1)\bigr\},\\
  n=2m+4 .
\end{gathered}
\]
As $m$ is even, $4\mid n$.

Take
\[
  S=\bigl\{[0,1],[0,3],\ldots,[0,m-1]\bigr\}\cup\bigl\{[1,0]\bigr\},
  \qquad |S|=\frac m2+1=\frac n4 .
\]
The cell $[0,2i-1]$ has vertex set $\{0,1\}\times\{2i-1,2i\}$, and as $i$ runs
over $1,\ldots,m/2$ these sets partition $\{0,1\}\times\{1,\ldots,m\}$; so
every vertex of the column except $(0,0)$ and $(1,0)$ is covered exactly once.
The cell $[1,0]$ covers $(1,0),(2,0),(1,1),(2,1)$. Of these, $(2,0)$ and
$(2,1)$ were uncovered and become covered once, $(1,0)$ was uncovered and
becomes covered once, and $(1,1)$ was already covered once and becomes covered
twice. The vertex $(0,0)$ is covered by no member of $S$. Thus $S$ is a
single-defect cover with $p=(0,0)$ and $v=(1,1)$, and
Lemmas~\ref{lem:cover} and \ref{lem:window} give both displayed assertions of
the theorem.

That $\Spec_{C_4}(L_m)$ is nonempty is the main theorem of \cite{CLE}, and
\eqref{eq:sym} makes it symmetric about $2(n+1)$; so if $s_0$ is any element
then so is $4(n+1)-s_0$, and the two lie strictly on opposite sides of the
centre, which is absent. Hence $\min<2(n+1)<\max$ with the centre missing, and
the set is not an interval.
\end{proof}

Parity is forced rather than chosen: $|S|=n/4$ must be an integer, and
$n=2m+4$ is divisible by $4$ exactly when $m$ is even. For odd $m$ the
argument above says nothing whatever, and we make no claim there.

\medskip
\noindent\textbf{Not claimed.}
That all six window values are actually attained is proved here only for
$L_2$, by the labelings of Table~\ref{tab:L2}; for larger even $m$ we assert
only nonemptiness, which is \cite{CLE}, together with the six-value window.
Nothing here shows that $L_2$ is a polyomino of least order with a
non-interval spectrum, nor that non-interval spectra arise only from
single-defect covers, nor that the gap must always be a single value at the
centre. Question~3 of \cite[Section~5]{CLE}, which asks for a
characterisation of the polyominoes $P$ with $2(n+1)\in\Spec_{C_4}(P)$,
remains open: Lemma~\ref{lem:cover}(iii) gives a sufficient condition for
failure, not a necessary one.

\section{Verification}\label{sec:verify}

Everything above can be checked by hand: one adds four numbers per cell in
Table~\ref{tab:L2} and checks the covering multiplicities of
$\{[0,1],[1,0]\}$ and of the family cover against the vertex lists.

A program \texttt{verify.py} accompanies this paper, together with a recorded
run \texttt{verify.output.txt}. It is Python~3.9+, standard library only, with
no external data file: the cell lists, covering sets and labelings are
literals in the source, transcribed from the paper. It re-derives $V(P)$ and
$n$ from the cell lists, verifies that each printed labeling is a bijection
onto $[n]$ and that every cell sums to the declared constant, recomputes the
covering multiplicities and the resulting values of $\delta$, checks the
arithmetic identity $(n/4)\cdot2(n+1)=n(n+1)/2$ that forces the collision at
the centre, confirms that the realised spectrum coincides with the six-value
window of Lemma~\ref{lem:window} and is symmetric under \eqref{eq:sym}, and
re-checks the covering rule of Theorem~\ref{thm:family} for all one hundred
even $m$ with $2\le m\le200$. Its recorded run also checks labelings and
covering profiles of further shapes that this note does not discuss. All
arithmetic is exact integer arithmetic; there is no search for labelings
anywhere, no floating point and no randomness, and the run takes well under a
second. It reports $99$ checks, all passing, and prints its own list of what
it does not cover.

\begin{thebibliography}{9}

\bibitem{CLE}
P.~Chalise, R.~M. Low, and A.~Eisenkolb-Vaithyanathan,
\emph{All polyominoes are $C_4$-face-magic},
preprint, submitted 9 August 2026.
\href{https://arxiv.org/abs/2608.08914v1}{arXiv:2608.08914v1}.
Section, definition and question numbers are quoted from v1 of that e-print,
the only version available to us; Question~2 is reproduced verbatim in
Section~1 above, so nothing here depends on the numbering.

\bibitem{Curran}
M.~J. Curran,
\emph{$C_4$-face-magic labelings on odd order projective grid graphs},
Australas. J. Combin. \textbf{83} (2022), no.~3, 361--396.

\bibitem{SLL}
W.~C. Shiu, R.~M. Low, and A.~K. Liu,
\emph{Face-magic labelings of polygonal graphs},
Theory Appl. Graphs \textbf{11} (2024), no.~1, Art.~7.
\href{https://doi.org/10.20429/tag.2024.110107}{doi:10.20429/tag.2024.110107}.

\end{thebibliography}

\end{document}
